% -----------------------------------------------------------------------------------
\section{The witness problems revisited, and counterpossibles}\label{sec:counterpossibles}

\subsection{Verdicts of the candidate semantics}

The witness problems of Section~\ref{sec:witnesses} can now be given as programs and evaluated
under each candidate semantics. Let $0$ stand for one-boxing or cooperating and $1$ for two-boxing or
defecting, and use the Newcomb and twin payoffs $u(a,b)=2(1-b)+a$, where $b$ is the predictor's
output or the twin's action. With $t=T(d)$ the reference agent's output, write the vectors of values
in the action order $(0,1)$ as
\[
  P_t=\bigl(2(1-t),\,2(1-t)+1\bigr),\quad S=(2,1),\quad D_t=(t,\,1-t),\quad Z=(0,0),\quad H=(5,10).
\]
$P_t$ freezes the predictor; $S$ transports a shared copy; $D_t$ and $Z$ are the frozen and transported
vectors of the mismatch diagonal $\one[a\ne b]$ of Theorem~\ref{thm:fixedcopy}; and $H$ gives the
payoffs of five-and-ten. In Newcomb and the twin the dependence is a designated oracle call; in the
diagonal the copy of the agent is a fixed component.

\begin{center}
\small
\renewcommand{\arraystretch}{1.2}
\begin{tabular}{@{}lcccc@{}}
  \toprule
  semantics & Newcomb & twin & five-and-ten & fixed-copy diagonal \\
  \midrule
  port only & $P_t$: two-box & $P_t$: defect & $H$: ten & $D_t$: opposite of $t$ \\
  designated slots & $S$: one-box & $S$: cooperate & $H$: ten & $D_t$: opposite of $t$ \\
  slots and literal copies & $S$ & $S$ & $H$ & $Z$: tie \\
  matching & $S$ & $S$ & $H$ & $Z$: tie \\
  supplied graph & $S$ or $P_t$ & $S$ or $P_t$ & $H$ & $Z$ or $D_t$ \\
  guarded, arithmetic & actual only & actual only & actual only & actual only \\
  logical induction & none & none & none & none \\
  \bottomrule
\end{tabular}
\end{center}

The first two rows are $\kappa_{\mathrm{port}}$ and $\kappa_{\mathrm{slot}}$, that is, $O_{\mathrm{const}}$
and $O_{\mathrm{oracle}}$; matching is local or global; logical induction is at an unspecified stage.
For the dependency graph the first entry applies when the graph links the predictor or copy to the
decision node and the second when it treats it as fixed. \enquote{Actual only} means that every
non-actual action has no value (Proposition~\ref{prop:guardedarith}), so there is no recommendation
under \ax{CF10}. For logical induction, nothing in the criterion fixes a numerical verdict at a given
stage; the all-zero-price stage of Section~\ref{sec:candidates} assigns the normalised upper bound to
both actions, giving $(3,3)$, $(3,3)$, $(10,10)$, and $(1,1)$.

The table makes visible what Section~\ref{sec:witnesses} took as given. In Newcomb's problem and the
twin prisoner's dilemma the conflict between \cri{1} and \cri{2} is the difference between the second
row and the first: whether the predictor's output, or the twin's action, is transported with the
agent's output or frozen at its baseline. A semantics that transports it one-boxes and cooperates;
one that freezes it two-boxes and defects. The model of Section~\ref{sec:model} made this choice by
stipulating $\mu_\pi$; a program does not make it, and the rows below the second show the ways an agent
may try to make it from source. On the fixed-copy diagonal, conversely, transporting the copy turns a
lost problem into a tie, which is why semantics that recognise copies escape
Theorem~\ref{thm:fixedcopy} and why Theorem~\ref{thm:complement} needs a copy that is transformed rather
than equal.

\subsection{Counterpossibles}

Every non-actual output of a deterministic computation is, relative to the computation's complete
description, an impossibility. Logical counterfactuals are therefore \emph{counterpossibles},
conditionals with impossible antecedents, and the philosophical literature on them bears on the
results of Section~\ref{sec:operators}.

On the orthodox semantics of \textcite{lewis1973counterfactuals}, a counterfactual with an impossible
antecedent is vacuously true. Vacuism is the counterfactual analogue of material implication in
Section~\ref{sec:logical}: it assigns every consequent to every non-actual output, and so provides a
relation, not a kernel. Non-vacuists answer with impossible worlds, in which impossible antecedents
hold without everything holding \autocite{nolan1997,brogaardsalerno2013}. \textcite{williamson2015,williamson2018}
defends vacuism, and one of his arguments has the structure of Theorem~\ref{thm:transport}. Suppose the
intension of $A\mathrel{\Prov\!\!\to}B$ is a function of the intensions of $A$ and of $A\wedge B$. If $A$ is
impossible, both intensions are empty for every $B$, so all counterfactuals with antecedent $A$ have the
same intension; since $A\mathrel{\Prov\!\!\to}A$ is true everywhere, every counterpossible is true.
\textcite{bertoetal2018} reply by admitting impossible worlds whose valuations discriminate more finely,
and discuss the conflict between reflexivity, closure under classical consequence, and non-trivial
treatment of contradictions, together with failures of substitution inside counterfactual contexts.

The analogy is exact at one level: an identification appropriate to actual or possible behaviour need
not preserve the distinctions needed under counterpossible suppositions. Williamson's argument
identifies antecedents by their possible-worlds intension, which is empty for every impossibility;
grounded extensionality identifies components by their instantiated behaviour, which is the same for
the identity context and the constant context at the baseline. In both cases the identification
collapses exactly the distinctions that non-vacuous counterfactuals need.

Impossible worlds can represent every kernel of Section~\ref{sec:operators}. Given a baseline-correct
kernel on finite environments, retain an actual world with the actual action and payoff; for every
defined alternative $a$ introduce a world recording action $a$ and payoff $\kappa(E,b,a)$, selected when
the antecedent \enquote{the output is $a$} is evaluated; and for an undefined alternative select several
worlds with incompatible payoffs, so that no value is determined. The selected worlds are impossible
relative to the complete description of the computation. Retaining the origin of each world and
overwriting its target realises \ax{CF7}. The construction represents $\kappa_{\mathrm{slot}}$,
$\kappa_{\mathrm{port}}$, $\kappa_{\mathrm{base}}$, and $\kappa_{\mathrm{flat}}$ alike, and therefore
derives none of them: impossible worlds can record unattainable values, which \ax{CF9} separately
excludes, and the selection relation encodes a counterfactual policy chosen in advance. Impossible-world
semantics makes room to preserve distinctions; it does not recover dependence information that a
closed program's actual behaviour has lost.

Two further precedents are close. \textcite{lewis1976} proved a triviality theorem about representing
conditional probabilities as probabilities of fixed conditional propositions; it does not prohibit
numerical conditional expectations or any kernel of Section~\ref{sec:operators}, but it is the classical
example of a semantic identification that cannot respect the required distinctions.
\textcite{taylor2020} discusses changing the agent's source code together with its hypothetical policy,
and the underdetermination of such changes, as an alternative to logical counterfactuals; it is closely
related to the slot semantics but contains no axiomatic impossibility. The discussion of embedded
agency by \textcite{demskigarrabrant2019} makes the point that an externally supplied agent slot already
contributes disputed structure, without a theorem about numerical kernels.

\subsection{What an optimality theorem requires}

The results of this part combine into a list of the assumptions that any optimality theorem for an
agent reading descriptions must state, because each of them changes some answer above:
\begin{enumerate}[label=(\roman*)]
  \item the atomic interfaces at which the agent is consulted;
  \item which contexts are intervention-total;
  \item whether implementations are compared on preserved inputs;
  \item whether fixed algorithms are admitted as constants;
  \item whether agents must halt globally or only on a promise class;
  \item the substitution benchmark by which alternatives are realised;
  \item how undefined alternatives are treated in ratifiability.
\end{enumerate}
They are mathematical assumptions, not implementation details that may be dropped once the theorem is
stated. Supplying them is supplying what Section~\ref{sec:model} supplied with $\mu_\pi$.
